\chapter{Zur kirchlichen Chronik}
\quellseite{6}

\section{America}

Unsere Brüder in Brasilien machen, ganz ähnlich wie die Väter und Gründer unserer Synode im Jahre 1839, jetzt schon den Anfang zur Gründung einer kleinen Lehranstalt. In der neuesten Nummer ihres vor Kurzem ins Leben gerufenen Blattes heißt es darüber: „Ein Institut zur Ausbildung von Lehrern und Predigern.

\quellseite{7}

Der obige Titel bezeichnet ein Werk, das erst im Werden begriffen ist, wodurch aber, wie wir hoffen, Ströme geistlichen Segens über unsere Kirche in diesem Lande fließen werden. Zweck und Ziel dieses Unternehmens ist: begabten jungen Leuten so viel theologischen und wissenschaftlichen Unterricht zu geben, daß sie als Lehrer oder Prediger der Kirche in ihrem eigenen Lande dienen können. Ein Zweifaches stand dabei vor Augen: zunächst die hohen Reisekosten für auswärtige Arbeiter in Zukunft zu sparen, sodann (und das ist der Hauptpunkt), die hiesige Kirche mit hier geborenen und aufgewachsenen Dienern zu versorgen, daß also die Kirche, so zu sagen, aus sich selbst wachse und erstarke. Dann erst, wenn die hiesige Kirche selbst ihre Lehrer und Prediger ausbildet, kann sie hoffen, einmal selbständig und vom Ausland unabhängig zu werden. Ein gesundes Kind will doch zuletzt selbst stehen und gehen lernen. Nach den hiesigen Verhältnissen, besonders auf den Colonien, wird ein zweifacher Cursus eingerichtet werden, einer, in dem Jünglinge in etwa vier Jahren zu Schullehrern herangebildet werden. Solche könnten eine Stelle übernehmen, wo ein Pfarrer aus der Umgegend die Kirchendienste versieht. In diesem Cursus würde eine gründliche Kenntniß der Sachen, die in der Schule Verwendung finden, vermittelt werden, sowie eine genügende Kenntniß der deutschen und portugiesischen Sprache, und nothdürftige Einführung in die höheren Rechenarten, Geschichts- und Erdkunde und ähnliche Gegenstände. Der zweite Cursus zur Ausbildung von Predigern wird naturgemäß längere Zeit umfassen und demgemäß der Unterricht auch weiter gehen. In beiden Abtheilungen wird vor allem auf Geistes- und Charakterbildung geachtet werden, da diese doch nöthiger und nützlicher ist als Verstandesbildung. Vorläufig ist der Ort für dieses Institut innerhalb der Gemeinde Bom Jesus auf der Colonie Sao Lourenco (P. F. Hartmeister). Ein bescheidenes Local zur Aufnahme etlicher Studenten wird hergerichtet. Sollte sich im Laufe der Zeit ein günstigerer Ort zeigen, so mag es ja dahin auf die Dauer verlegt werden. Wie gesagt, das Werk ist erst im Werden begriffen. Vier Schüler sind schon gemeldet und warten auf die Eröffnung, die, so Gott will, im October stattfinden soll. Wer nun ein Herz für die Sache hat, nachdem er ihre Nothwendigkeit und Nützlichkeit eingesehen hat, dem sei sie zu warmer Fürbitte herzlich empfohlen.

L. F.

In den stockkatholischen südamericanischen Ländern, wie Peru, Ecuador und Bolivia, waren bisher alle öffentlichen gottesdienstlichen Versammlungen, die nicht von einem katholischen Priester geleitet wurden, vom Staate verboten. Nun ist es den vereinten Bemühungen amerikanischer Kirchen gelungen, die Regierungen der Vereinigten Staaten, Englands und Deutschlands zu bewegen, durch ihre Gesandten in Peru bei der dortigen Regierung vorstellig zu werden, daß die Abhaltung protestantischer Gottesdienste freigegeben werde. Daraufhin ist auch eine Vorlage, die solche Gottesdienste gestattet, in der gesetzgebenden Versammlung eingereicht worden, und die Hoffnung wird ausgesprochen, daß auch in den andern Ländern den Protestanten mehr Religionsfreiheit gewährt werden wird. Die Presbyterianer, Methodisten, Episkopalen und Congregationalisten haben in Peru Missionen, die unter dem protestantenfeindlichen Gesetze gelitten haben. Daß es dort auch gar manche zerstreute Lutheraner gibt, die von Deutschland aus nur nothdürftig kirchlich versorgt werden, haben wir schon früher einmal mitgetheilt. Die Hauptstation der Lutheraner ist Callao, in der Nähe der Hauptstadt Lima.

L. F.

\section{Ausland}

\subsection{Frauenemancipation}

Der deutsche Verein für Frauenstimmrecht hat, wie der „Lutherische Herold“ mittheilt, an die bekanntesten lutherischen, reformirten und katholischen Professoren der Theologie ein Schreiben mit folgenden Anfragen gerichtet: 1. ob sich aus den

Reden Jesu ein directes Verbot der vollen Gleichberechtigung von Mann und Frau ableiten läßt, und 2. ob in den Verfassungen der christlichen Urgemeinden eine Spur eines Unterschiedes in den Rechten männlicher und weiblicher Gemeindemitglieder nachweisbar ist. — Die Fragen sind recht vorsichtig gestellt und so formulirt, daß die Antworten womöglich günstig für die Frauen lauten sollen. Und doch ist die ganze Anfrage höchst überflüssig. Wenn diese Frauen wirklich wissen wollen, was die Schrift von der modernen Frauenbewegung urtheilt, so brauchen sie nur die Bibel aufzuschlagen und zu lesen, was der Apostel Jesu Christi, St. Paulus, darüber zu sagen hat. Der beruft sich wiederholt auf die alte, unverbrüchliche Ordnung Gottes: „Dein Wille soll deinem Mann unterworfen sein, und er soll dein Herr sein“, 1 Mos. 3, 16. Darum bezeichnet er den Mann als des Weibes Haupt, 1 Cor. 11, 3. Eph. 5, 22. ff. Col. 3, 18., und weiß das Weib aus dem öffentlichen Leben und aus der öffentlichen Thätigkeit in den Bereich des Hauses und der Familie, 1 Cor. 14, 34. f. 1 Tim. 2, 11. ff. Wo man sich freilich über diese Worte des Heiligen Geistes hinweggesetzt und darüber spöttelt, da würde man auch einem besonderen Worte Christi kein Gehör schenken.

L. F.

\subsection{Großbritannien}

steht an der Spitze aller Länder, was die für die Heidenmission aufgebrachten Gelder anlangt. Die Jahreseinnahmen seiner zehn größten Missionsgesellschaften betrugen im letzten Jahre in runden Zahlen zusammen \$7,000,000. Oben steht die “Church Mission Society” mit rund \$1,750,000; das ist mehr, als alle deutschen Missionsgesellschaften zusammen vereinnahmen. Ihr folgen die Mission der schottischen Freikirche, die Londoner Mission, die der wesleyanischen Methodisten, die Gesellschaft zur Ausbreitung des Evangeliums, die Mission der Baptisten, die der schottischen Staatskirche und die China-Inland-Mission. Jede dieser Gesellschaften hat eine Jahreseinnahme von mehr als \$250,000. Dazu kommen dann noch zwei Hülfsgesellschaften der Heidenmission, die in dieser Zusammenstellung nicht mit eingriffen sind, nämlich die Britische und Ausländische Bibelgesellschaft mit einer Jahreseinnahme von rund \$1,200,000 und die Tractatgesellschaft mit \$500,000 Einnahme. Nach England folgen in den Jahreseinnahmen die Missionsgesellschaften der Vereinigten Staaten, dann der Reihe nach Deutschland, die skandinavischen Länder, Canada, Frankreich, Holland. — Es sind in der That große Summen, die für Heidenmission aufgebracht werden. Und doch stehen sie in keinem Verhältniß zu dem irdischen Reichthum der sichtbaren Christenheit.

L. F.

\subsection{Der neue Papst Pius X.}

hat die schönen Hoffnungen, die manche auf ihn gesetzt hatten, die eben das Papstthum nicht kennen, sehr schnell zu Schanden gemacht. Sein Vorgänger Leo XIII. war bekanntlich ein schlauer Politiker, der unablässig darauf bedacht war, die äußerliche, politische Macht des Papstthums zu befestigen. Pius dagegen wurde als ein „stiller, bescheidener, frommer, religiöser“ Mann hingestellt, der sich wenig um Politik und um die Machtstellung des Papstthums bekümmern, vielmehr allem Gewicht auf Religion und Frömmigkeit legen werde. Insbesondere hatte man gehofft, daß er die gespannten Beziehungen seiner Vorgänger zum italienischen Königshause nicht aufrecht erhalten werde. Auch das Blatt der deutschen Reformkatholiken, das „Zwanzigste Jahrhundert“, schrieb: „Auf dem Thron wird Pius X. — dessen können wir sicher sein — kein anderer sein, als er als einfacher Priester und Bischof war. Er wird demüthig und bescheiden sein, sich nicht um irdischen Pomp, sondern um Ehre und Herrlichkeit im Jenseits kümmern: er wird der religiöse Papst sein, den die Völker ersehnt haben in unsern Tagen, wo sie den politisirenden Priester nicht versehen und nicht ertragen wollen.“ Aber in seiner ersten amtlichen Ansprache an die versammelten Cardinäle am 9. November schlug Pius einen ganz andern Ton an und gab deutlich zu erkennen,
\quellseite{8}
daß er genau den Weg seiner Vorgänger gehen werde. Zuerst kam er auf seine Wahl zum Papste zurück und sagte in heuchlerischer, scheinheiliger Weise, daß er vergeblich versucht habe, das Amt abzuweisen; schließlich habe er es als Gottes Willen erkannt, daß er es annehmen solle. Wenn der Papst auch nur einen Augenblick nach Gottes Willen fragen wollte, dann würde er sofort das ganze Papstthum aufheben müssen. Denn Christus will keinen solchen Herrn und angeblichen Stellvertreter seiner selbst in der Kirche haben, sondern spricht: „Einer ist euer Meister, Christus; ihr aber seid alle Brüder“, Matth. 23, 8. 10. — Dann kommt der Papst auf sein päpstliches Programm, auf die von ihm zu befolgende Politik, zu sprechen und sagt da nach dem Bericht der katholischen „Kölnischen Volkszeitung“ wörtlich Folgendes: „Da es nöthig ist und auch der christliche Staat das größte Interesse daran hat, daß der Papst in der Leitung der Kirche frei sei und als frei erkannt werde und keiner Macht unterthan sei, so zwingt uns die gewissenhafte Auffassung unseres Amtes und der von uns geleitete Eid, das in diesem Punkte der Kirche zugefügte schwere Unrecht zu beklagen.“ Es ist aber nicht wahr, daß der Papst bloß Freiheit in der Leitung der Kirche haben will. Er will vielmehr zugleich ein weltlicher Herrscher sein und in weltlichen Händeln mitthun, während doch Christus gesagt hat: „Mein Reich ist nicht von dieser Welt“, Joh. 18, 36. Deshalb beklagt sich auch der Papst über das „der Kirche zugefügte schwere Unrecht“. Er meint damit die Wegnahme des Kirchenstaates, die Aufhebung der weltlichen Herrschaft des Papstes durch die italienische Regierung. Pius sagt dann weiter: „Man muß sich darüber verwundern, daß es Leute gibt, die in ihrer Neugierde sich mit Vermuthungen abquälen über unser Regierungsprogramm. Als ob das einer Untersuchung bedürfte, oder als ob es nicht ganz offen kundig sei, daß wir denselben Weg gehen wollen und müssen, den unsere Vorgänger bis heute eingehalten haben .... Wenn man Anstoß daran nimmt, daß wir uns um Politik kümmern, so muß doch jeder billig Denkende einsehen, daß der Papst von seiner lehramtlichen Thätigkeit auf Glaubens- und Sittengebiet Fragen politischer Art nicht trennen kann. Ueberdies muß er als Oberhaupt und Lenker einer vollkommenen, aus Menschen sich zusammen setzenden und unter Menschen bestehenden Gesellschaft, wie es die Kirche ist, doch zweifellos den Wunsch haben, daß zwischen ihm und den Staatsoberhäuptern und Staatslenkern Beziehungen bestehen, wenn er für die Sicherheit und Freiheit der Katholiken in allen Welttheilen gesorgt wissen will.“ Die Welt hat also von dem jetzigen Papste ganz dasselbe zu erwarten wie von seinen Vorgängern. Er wird dieselben unverschämten Ansprüche in geistlichen und weltlichen Dingen erheben, dieselben Verdammungsurtheile aussprechen über diejenigen, die ihm nicht gehorchen, dieselben Schmähungen und Lästerungen ausstoßen gegen unsern Herrn und Heiland Jesum Christum und sein theueres Verdienst, dieselben greulichen Mißbräuche und seelenverderblichen Irrlehren festhalten. Das kann uns nicht Wunder nehmen. Papstthum bleibt Papstthum, wer auch immer gerade Papst ist. Das Papstthum ist eben nicht etwas, was von den einzelnen Inhabern des päpstlichen Stuhles abhängt. Es ist ein System, eine Einrichtung, die, wie schon Luther erkannte und sagte, nur so gebessert und reformirt werden kann, daß sie ganz und gar abgeschafft wird. Das Papstthum ist eine Einrichtung ganz des Satans, und erst der Herr wird ihm ein Ende machen durch die Erscheinung seiner Zukunft, 2 Thess. 2, 8. f.

Zur Zeit der französischen Revolution, am 10. November 1793, wurde bekanntlich in Paris die christliche Religion ganz förmlich abgeschafft und die Vernunft als Göttin erhoben, anerkannt und gefeiert. Eine Schauspielerin, gekleidet in die Farben der Revolution, wurde im Triumphwagen durch die Straßen geführt, in der Pariser Hauptkirche Notre dame auf den Hauptaltar gehoben und mit Ver-

L. F.

neigungen, Räucherungen, Ansprachen und Gesängen als „Göttin der Vernunft“ verehrt. In den andern Pariser Kirchen und in mehreren tausend Kirchen des Landes wurde diese Feier nachgeahmt. Und jetzt kommt aus Paris die Nachricht, daß die dortigen Atheisten am 8. November auf Anregen des ehemaligen katholischen Priesters Charbonnel zur Erinnerung an jenen Vorgang eine neue Feier der Vernunft in einem großen Saale abgehalten haben. Eine Anzahl Reden, mit Gedichten und Gesängen untermischt, wurden der einige tausend Köpfe zählenden Versammlung geboten. Einen Gott gebe es nicht, die christliche Kirche habe ihren Einfluß verloren, kein Gebildeter glaube mehr, was es nicht mit der Vernunft beareifen könne — in solchen und ähnlichen Redewendungen wurde den ganzen Abend gesprochen. — „Die Thoren sprechen in ihrem Herzen: Es ist kein Gott. Sie taugen nichts, und sind ein Greuel mit ihrem Wesen“, Ps. 14, 1. „Der im Himmel wohnet, lachet ihrer, und der Herr spottet ihrer“, Ps. 2, 4.

L. F.
